\section{Simplified waste-generation schemes}\label{sec:sp:schemes}

The general scheme of Section~\ref{sec:sp:setup:materialaccounting} lets waste arise from every use of
the final good and, separately, from two physical flows measured in tonnes. Much of it can be switched
off by setting parameters to zero. The switches below are used to build up, test and decompose the
model; the useful question about each is not whether it is realistic but whether it changes the
allocation or only the reported mass flows. That question has different answers under the two versions,
and the difference is the sharpest practical reason to carry both.

Throughout, references to Version B are to Section~\ref{sec:sp:B} and references to Version A to
Section~\ref{sec:sp:A}. Recall the two objects the classification turns on: the relative intensities
$\omega^j$ scaling the streams tied to embodied material in $Y$, and the composition wedge
$\Delta=\Omega\chi\left(p^W-p^M\right)$ of~\eqref{eq:sp:A:zeta}, which is identically zero under
Version B.

\begin{enumerate}[label=\textbf{(W\arabic*)},ref=W\arabic*,leftmargin=3.2em,itemsep=0.35em]

\item\label{sw:sp:nature} \textbf{No nature-attributable waste:} $\Omega^{N,S}=\Omega^{D,S}=0$, so
$\mathcal N\equiv0$. Overburden and gangue are ignored, and every tonne of waste has passed through
production. The collected ledger~\eqref{eq:sp:setup:ledger-collected} becomes
$W\left(1-a\varpi\right)=N-\dot M^K$; the extraction and exploration conditions lose their $p^W$ terms,
giving $B^{\mathrm B}=C^N_N+p^S+d^Wp^P$ and $p^S+p^X=C^D_D$ under Version B, and the corresponding wedged versions
under Version A; and the waste price becomes explicit rather than implicit. Under Version B,
\begin{align}
p^W &= c^W\!\left(\varpi\right)+\left(1-\varpi+d^W\varpi\right)p^P
-\varpi\alpha\left(C^N_N+p^S+d^Wp^P\right),
\label{eq:sp:schemes:pW-simple}
\end{align}
which is~\eqref{eq:sp:B:foc:W-solved} with the $\Omega^{N,S}\varpi\alpha$ term gone. Under Version A the
same simplification removes the same term from~\eqref{eq:sp:A:foc:W-solved}, but $p^W$ remains implicit
there through $\Delta$. \emph{Affects the allocation under both versions.}

\item\label{sw:sp:omegaY} \textbf{Waste from output only:} $\omega^Y=1$ and
$\omega^C=\omega^{N}=\omega^{D}=\omega^W=0$, the textbook $W\propto Y$ scheme. Under Version A this
leaves a single wedge, on output and on the value of materials, and removes it from consumption,
extraction, exploration and treatment.

\item\label{sw:sp:omegaYC} \textbf{Waste from output and consumption only:}
$\omega^{N}=\omega^{D}=\omega^W=0$, leaving $\omega^Y,\omega^C$ free. Under Version A this is the
smallest scheme that retains a consumption wedge, and therefore the smallest one in which the
version choice has any bearing on the consumption--saving margin.

\item\label{sw:sp:uniform} \textbf{Uniform intensities:} $\omega^j=1$ for every $j$, so that
$\mathcal D$ is total final-good use and $\Omega$ is a single economy-wide intensity. Under Version A
every wedge is then the same scalar $\Delta$, applied with a plus sign to spending and a minus sign to
output.

\item\label{sw:sp:phiI} \textbf{Capital embodies nothing:} $\Phi^I=0$ (Version B) or $\phi^I=0$
(Version A). Then $M^K\equiv0$, the state and its price $p^M$ drop out, $q=1$, the mass
constraint~\eqref{eq:sp:B:mass} is vacuous, and by Remark~\ref{rem:sp:coincide} the two versions
coincide --- $\chi=0$ kills $\Delta$, so Version A's entire apparatus collapses. The ledger reduces to
$W=\left(N+\mathcal N\right)/\left(1-a\varpi\right)$ throughout, not only at $\dot M^K=0$. This is the
cleanest available stripped model and the natural first target for the numerical solver.
\emph{Affects the allocation under both versions.}

\item\label{sw:sp:explore} \textbf{No exploration:} $D\equiv0$ and $C^D\equiv0$. The states $X$ and $S$
reduce to a single depleting reserve, $\dot S=-N$, and $p^X$ drops out; $p^S$ becomes the pure
modified-Hotelling rent~\eqref{eq:sp:B:costate:S}. \emph{Affects the allocation under both versions},
and note that it does \emph{not} restore a steady state --- see~\ref{sw:sp:CDX}.

\item\label{sw:sp:CDX} \textbf{Exploration costs independent of past discovery:} $C^D_X=0$, so $X$
ceases to be a state and $p^X\equiv0$. The exploration condition becomes $p^S=C^D_D+\Omega^{D,S}p^W$
under Version B, and $p^S=\left(1+\Delta\omega^D\right)C^D_D+\Omega^{D,S}p^W$ under Version A. This is
the switch that repairs Propositions~\ref{prop:sp:B:nosteadystate}
and~\ref{prop:sp:A:nosteadystate}: with $X$ gone, $\dot S=D-N=0$ is compatible with $D=N>0$, and a
genuine interior steady state exists, characterised under Version B by the quasi-stationary
block~\eqref{eq:sp:B:quasistationary} together with
$N\left(1+\Omega^{N,S}+\Omega^{D,S}\right)=W\left(1-a\varpi\right)$ and $\Phi^K=\Phi^I$.
\emph{Affects the allocation under both versions.}

\item\label{sw:sp:nopoll} \textbf{No pollution channel in production:} $F_P=0$, so damages enter only
through $v(P)$, with $d^{\mathrm B}=v'(P)/u'(C)$ under Version B and $d=v'(P)/\lambda$ under Version A.
\emph{Affects the allocation under both versions}, though only through the level of $p^P$.

\end{enumerate}

\smalltitle{The classification.} Switches~\ref{sw:sp:omegaY}--\ref{sw:sp:uniform} are of a different
kind from the rest, and the difference is exactly
Proposition~\ref{prop:sp:B:irrelevance}.

\begin{proposition}[Waste-composition switches are allocationally inert under Version B]
\label{prop:sp:inert}
Under Version B, any restriction on the relative intensities $\left\{\omega^Y,\omega^C,\omega^{N},
\omega^{D},\omega^W\right\}$ --- including~\ref{sw:sp:omegaY}, \ref{sw:sp:omegaYC}
and~\ref{sw:sp:uniform} --- leaves the optimal paths of every state, every control and every shadow
price unchanged. It changes only the derived quantities $\Omega$ and $\Phi^Y$
in~\eqref{eq:sp:B:Omega}. Under Version A the same restrictions change every condition
in~\eqref{eq:sp:A:foc:C}--\eqref{eq:sp:A:costate:M}.
\end{proposition}

The first part is immediate from Proposition~\ref{prop:sp:B:irrelevance}: no $\omega^j$ appears anywhere
in Problem~\ref{prob:sp:B} or in its optimality conditions, so no restriction on them can move the
solution. The second is immediate from inspection of the factors $\left(1\pm\Delta\omega^j\right)$ in
Section~\ref{sec:sp:A:foc}, which are the only place the intensities enter and which are not identically
one. Table~\ref{tab:sp:switches} collects the classification.

\begin{table}[H]
\centering
\caption{Simplifications of the waste-generation scheme, and what each affects.}
\label{tab:sp:switches}
\setlength{\tabcolsep}{5pt}
\renewcommand{\arraystretch}{1.05}
\begin{tabular}{@{}p{0.30\textwidth}p{0.20\textwidth}p{0.20\textwidth}p{0.22\textwidth}@{}}
\toprule
Switch & Version B & Version A & States removed \\
\midrule
\ref{sw:sp:nature}\ \ no $\mathcal N$          & allocation & allocation & --- \\
\ref{sw:sp:omegaY}\ \ waste from $Y$ only      & reporting only & allocation & --- \\
\ref{sw:sp:omegaYC}\ \ waste from $Y,C$ only   & reporting only & allocation & --- \\
\ref{sw:sp:uniform}\ \ uniform $\omega^j$      & reporting only & allocation & --- \\
\ref{sw:sp:phiI}\ \ $\Phi^I=0$                 & allocation & allocation (collapses to A) & $M^K$ \\
\ref{sw:sp:explore}\ \ no exploration          & allocation & allocation & $X$ \\
\ref{sw:sp:CDX}\ \ $C^D_X=0$                   & allocation & allocation & $X$ \\
\ref{sw:sp:nopoll}\ \ $F_P=0$                  & allocation & allocation & --- \\
\bottomrule
\end{tabular}
\end{table}

Two readings of the table matter for what follows. First, the only simplifications that are free under
Version B are precisely the ones about the \emph{composition} of waste across final-good uses --- the
part of the accounting that is hardest to discipline empirically and that occupies most of
Section~\ref{sec:sp:setup:materialaccounting}. Under Version B that section can be read as bookkeeping;
under Version A it cannot. Second, none of the switches removes the ledger itself. Mass conservation is
not a modelling option in this framework, and no setting of the parameters recovers a specification in
which recycling avoids waste rather than deferring it.

\smalltitle{A note for the numerical implementation.} Switch~\ref{sw:sp:phiI} is the natural entry
point, since it collapses the two versions into one and removes a state; switch~\ref{sw:sp:CDX} is the
natural second step, since it restores a steady state to converge to. Together they give a model with
three states, a single interest rate, no complementary-slackness block, and no dependence on the
relative intensities at all --- which is to say, a model in which every one of the open items of
Appendix~\ref{sec:sp:extensions} except~\ref{op:sp:period} is either resolved or vacuous. Everything
beyond that point is added back one switch at a time, and the classification above says which
additions require the solver to change and which only change what is reported afterwards.
